\documentclass[11pt,a4paper]{article}
\usepackage[T1]{fontenc}    % Indispensable pour la coupure des mots en français
\usepackage[utf8]{inputenc} 
\usepackage{amsmath}
\usepackage{amssymb}
\usepackage{amsfonts}
\usepackage{stmaryrd}
\usepackage[margin=1cm]{geometry}
\usepackage{frenchmath}

\begin{document}

% --- TITRE DE L'EXERCICE (Style identique au titre 1.7 de l'image) ---
\begin{center}
    \large \textbf{1. Etude d'une matrice à diagonale nulle}
\end{center}
\vspace{0.2cm}

% --- ENCADRÉ DE L'ÉNONCÉ (Même style visuel avec source en bas à droite) ---
\noindent\framebox[\textwidth]{%
    \begin{minipage}{0.94\textwidth}
        \vspace{0.2cm}
        Soit un entier $n \geqslant 2$. On note $A$ la matrice de $\mathcal{M}_n(\mathbb{R})$ dont les coefficients diagonaux valent 0 et les autres coefficients valent 1.
        \vspace{0.2cm}
        
        \noindent $1.$ Calculer $A^2$. En déduire que $A$ est inversible et exprimer $A^{-1}$ ;
        
        \noindent $2.$ Déterminer les valeurs propres de $A$ et les espaces propres correspondants.
        \vspace{0.1cm}
        \begin{flushright}
            (\textbf{Mines-Ponts, PC})
        \end{flushright}
    \end{minipage}%
}
\vspace{0.4cm}

% --- DÉBUT DE LA SOLUTION (Introduction avec triangle \triangleright) ---
\noindent $\triangleright$ \textbf{Solution.} \\

% --- PREMIÈRE QUESTION (Introduite par une puce \bullet comme dans l'image) ---
\textbf{1.} Commençons par calculer $A^2$. Soit $(i,j) \in \llbracket 1,n \rrbracket^2$. On a :
\[
[A^2]_{i,j} = \sum_{k=1}^n A_{i,k} A_{k,j} = \sum_{\substack{k=1 \\ k \notin \{i,j\}}}^n A_{i,k} A_{k,j} = \sum_{\substack{k=1 \\ k \notin \{i,j\}}}^n 1 = \left| \llbracket 1,n \rrbracket \setminus \{i,j\} \right|,
\]
ce qui donne:
\[
[A^2]_{i,j} = \begin{cases} n-2 & \text{si } i \neq j \\ n-1 & \text{si } i = j \end{cases}
\]
On en déduit que 
\[
A^2 = (n-2)A + (n-1)I_n.
\]
On peut réécrire cette relation sous la forme $A \cdot \left( \frac{1}{n-1}(A - (n-2)I_n) \right) = I_n$. Ainsi, \fbox{la matrice $A$ est inversible} et, par unicité de l'inverse, on obtient :
\[
\boxed{A^{-1} = \frac{1}{n-1}(A - (n-2)I_n).}
\]

\vspace{0.4cm}

% --- DEUXIÈME QUESTION (Deuxième puce \bullet) ---
\textbf{2.} Déterminons les valeurs propres de $A$. L'équation aux valeurs propres $AX = \lambda X$, avec $\lambda \in \mathbb{C}$ et $X = (x_1, \dots, x_n)^T \in \mathbb{R}^n \setminus \{0\}$, équivaut au système suivant :
\[
\forall i \in \llbracket 1,n \rrbracket, \quad \sum_{\substack{j=1 \\ j \neq i}}^n x_j = \lambda x_i \quad (1)
\]
En sommant ces $n$ identités, il vient :
\[
(n-1) \left( \sum_{j=1}^n x_j \right) = \lambda \left( \sum_{i=1}^n x_i \right).
\]
Deux cas se présentent alors:
\begin{itemize}
    \item[\textbf{--}] Si $\sum_{i=1}^n x_i = 0$ : le vecteur $X$ appartient à l'hyperplan $H = \left\{ X = (x_1, \dots, x_n)^T \in \mathbb{R}^n \;\left|\; \sum_{i=1}^n x_i = 0 \right.\right\}$. On déduit alors de la relation (1) que :
    \[
    \sum_{j=1}^n x_j - x_i = \lambda x_i \quad \forall i \in \llbracket 1,n \rrbracket,
    \]
    soit $-x_i = \lambda x_i$, d'où $\lambda = -1$.
    \item[\textbf{--}] Si $\sum_{i=1}^n x_i \neq 0$ : on obtient directement $\lambda = n-1$.
\end{itemize}
On peut donc déduire que le spectre de $A$ est :
\[
\boxed{\mathrm{Sp}(A) = \{n-1, -1\},}
\]
et les espaces propres associés sont donnés par :
\[
\boxed{E_{n-1}(A) = \mathrm{Vect}\left((1, \dots, 1)^T\right) \quad \text{et} \quad E_{-1}(A) = H.}
\]

\end{document}